\documentclass[10pt,a4paper]{amsart}
\usepackage[margin=19mm]{geometry}
\usepackage{amsmath,amssymb,mathtools}
\usepackage[scr=rsfs,cal=euler]{mathalfa}
\usepackage{xfrac}
\usepackage[unicode,hidelinks,bookmarks=false]{hyperref}
\newcommand{\ie}{i.e.\ }
\newcommand{\cf}{cf.\ }
\newcommand{\resp}{respectively\ }
\newcommand{\Subsection}{\subsection}
\providecommand{\nobreakdash}{\nobreak\hbox{-}}
\setlength{\emergencystretch}{2em}
\multlinegap0pt
\usepackage{amssymb}
\usepackage{mathtools}
\usepackage{xcolor}
\newtheorem{lemma}{Lemma}
\newtheorem{proposition}{Proposition}
\usepackage{cleveref}
\crefname{lemma}{Lemma}{Lemmas}
\crefname{proposition}{Proposition}{Propositions}
\def\K{\mathcal K}
\def\Kb{\mathcal K_b}
\def\M{\mathcal M}
\newcommand{\divg}{\operatorname{div}}
\title{\vspace{-1cm} Quantitative Diffusive Limits for Singular Nonlocal Transport}
\author{Andrea Agazzi, Giuseppe Bruno, Federico Pasqualotto, Philippe Rigollet}
\date{Source: arXiv:2609.11837v1 (CC BY 4.0)}
\begin{document}
\maketitle

\noindent\emph{Source and licence.} \vspace{-1cm} Quantitative Diffusive Limits for Singular Nonlocal Transport. Version arXiv:2609.11837v1 (CC BY 4.0), distributed by arXiv under the Creative Commons Attribution 4.0 International licence, http://creativecommons.org/licenses/by/4.0/, as recorded in the arXiv licence metadata for this exact version and shown on the abstract page https://arxiv.org/abs/2609.11837v1. Full licence notice: \texttt{licenses/arxiv-2609.11837-CC-BY-4.0.txt}.\par\smallskip

\noindent\textit{Editorial scope.} Counted unit: the article's proposition prop:pde\_holder\_propagation (source0.tex lines 3314--3325). The article's complete original proof of it is delivered with it (source0.tex lines 3327--3577, one \texttt{proof} environment of 251 lines). No mathematics is written, repaired or re-derived: the statement and the proof are the article's own lines, in the article's own notation, and no retained line is substituted or re-flowed.  Retained uncounted as the source-local context the proof needs: lines 2627--2654; lines 2630--2632; lines 2964--2981; lines 3194--3200; lines 3250--3268.  Nothing was omitted from any retained span, so the package carries no omission record.  No retained line contains a citation, so the wrapper carries no bibliography and no bibliography entry.  No retained line contains a figure environment, an image-inclusion command or drawing code, so no image is delivered and the package declares no asset dependency.  Editorial formatting only, in addition: an AMS \texttt{amsart} preamble that restates the article's own declarations and macros used by the retained lines, namely the environments \texttt{lemma}, \texttt{proposition} on separate counters, together with the standard packages those lines need.  The retained environments print in this document as Lemma 1, Proposition 1 in order of appearance, whereas the article prints the same items with the numbering of its own full document.  The complete native source of the article as distributed in the arXiv e-print for this exact version is bundled unchanged as \texttt{sources/210055/source0.tex}.
\begin{lemma}[Fixed-\(b\) kernel estimates]\label{lem:wellposed_kernel_estimates}
Let \(M>0\). If \(\rho\) is a probability density on  \(\M\) with
\(\|\rho\|_{L^\infty}\le M\), and \(u=\Kb\rho\), then there exist $\kappa_b>0$ and $C_{b,M,\M}>0$ such that
\begin{equation}\label{eq:wellposed_fixed_b_basic}
\kappa_b\le u\le M,\qquad \|\nabla u\|_{L^\infty}\le C_{b,M,\M}.
\end{equation}

Furthermore, for every $x,y\in \M$ with $x\neq y$, let \(P_{y\to x}\) the parallel
transport along any minimizing geodesic from \(y\) to \(x\), then:
\[
|\nabla u(x)-P_{y\to x}\nabla u(y)|
\le C_{b,M,\M}\operatorname{dist}_{ \M}(x,y)\log(2+\operatorname{dist}_{ \M}(x,y)^{-1}).
\]
Consequently
\[
|V[\rho](x)-P_{y\to x}V[\rho](y)|
\le C_{b,M,\M}\operatorname{dist}_{ \M}(x,y)\log(2+\operatorname{dist}_{ \M}(x,y)^{-1}).
\]
If \(\rho\in C^{k,\alpha}( \M)\), \(k\ge0\), \(0<\alpha<1\), then for all $b>0$
\begin{equation}
\label{eq:schauder}
\|\K_b\rho\|_{C^{k+2,\alpha}}
\le C_{b,k,\alpha,\M}\|\rho\|_{C^{k,\alpha}},
\quad
\|V[\rho]\|_{C^{k+1,\alpha}}
\le C_{b,k,\alpha,\M}\bigl(1+\|\rho\|_{C^{k,\alpha}}\bigr)^{k+2}.
\end{equation}
\end{lemma}

\begin{equation}
\kappa_b\le u\le M,\qquad \|\nabla u\|_{L^\infty}\le C_{b,M,\M}.
\end{equation}

\begin{lemma}[Logarithmic bound]\label{lem:wellposed_log_lipschitz}
Let \(0<\alpha<1\). If \(f\in C^{0,\alpha}( \M)\), then
\begin{equation}
\label{eq:log_bound_hess}
\|\nabla^2\K_b f\|_{L^\infty}
\le
C_{b,\alpha,\M}\|f\|_{L^\infty}
\left(1+\log\left(1+\frac{[f]_{C^{0,\alpha}}}{\|f\|_{L^\infty}}\right)\right),
\end{equation}
with the convention that the right-hand side is zero when \(f=0\). Hence, if
\(\rho\in C^{0,\alpha}( \M)\) is a probability density with
\(\|\rho\|_{L^\infty}\le M\), then
\[
\|\nabla V[\rho]\|_{L^\infty}
\le
C_{b,M,\alpha,\M}\left(1+\log\bigl(2+\|\rho\|_{C^{0,\alpha}}\bigr)\right).
\]
\end{lemma}

\begin{equation}\label{eq:short_chizat_G}
\Delta\log u
=
\frac{\Delta u}{u}-|\nabla\log u|^2
=
\frac{u-\rho}{b^2u}-|\nabla\log u|^2 .
\end{equation}

\begin{proposition}[Local bounded well-posedness and continuation criterion]\label{prop:pde_local_bounded_wellposedness}
Let \(b>0\) and let \(\mu_{\mathrm{in}}\) be a bounded probability density on
 \(\M\).
Then there are \(T_{\max}\in(0,\infty]\) and a unique maximal bounded
solution of
\[
\partial_t\mu+\divg(\mu V[\mu])=0,\qquad V[\mu]=-\nabla\log \mathcal K_b\mu,
\]
with initial datum \(\mu_{\mathrm{in}}\), satisfying
\[
\mu\in L^\infty_{\rm loc}([0,T_{\max});L^\infty( \M)).
\]
It preserves mass, and the following continuation criterion holds:
\[
T_{\max}<\infty
\quad\Longrightarrow\quad
\limsup_{t\uparrow T_{\max}}\|\mu_t\|_{L^\infty( \M)}=+\infty .
\]
\end{proposition}

\begin{proposition}[Local propagation of Hölder regularity]\label{prop:pde_holder_propagation}
Let \(0<\alpha<1\), \(k\in\mathbb N\cup\{0\}\), and let
    \(\mu_{\mathrm{in}}\in C^{k,\alpha}( \M)\) be a bounded probability density.
Let \(\mu\) be the maximal bounded solution from
\Cref{prop:pde_local_bounded_wellposedness}, with maximal time
\(T_{\max}\).  Then
\[
\mu\in L^\infty_{\rm loc}([0,T_{\max});C^{k,\alpha}( \M)).
\]
If \(\mu_{\mathrm{in}}\in C^\infty\), then \(\mu\) is smooth in space and time on
\([0,T_{\max})\times \M\).
\end{proposition}

\begin{proof}
Fix \(T\in(0,T_{\max})\).  The proof proceeds in three stages: short-time Hölder regularity,
continuation, and a higher-order bootstrap.  
Set
\[
M_T:=\sup_{t\in[0,T]}\|\mu_t\|_{L^\infty}<\infty .
\]
Throughout the proof, \(C\) denotes a positive, finite constant whose value may
change from line to line.  In Steps 1--2 it depends only on
 \(\M\), \(b\), \(\alpha\), \(T\), \(M_T\), and
\(\|\mu_{\mathrm{in}}\|_{C^{0,\alpha}}\), while in Step 3 it may also depend on the
lower-order Hölder norms already controlled in the bootstrap.
Our equation can be written in an equivalent way as
\[
\partial_t\mu_t+V[\mu_t]\cdot\nabla\mu_t=\mu_t\Delta\log u_t.
\]

\smallskip
\noindent\textbf{Step 1.}
We first prove that \(\mu\in L^\infty([0,T_1];C^{0,\alpha})\) for some
\(0<T_1\le T\).  Consider the Picard sequence
\[
\mu^0_t=\mu_{\mathrm{in}},\qquad
\mu^{n+1}_t=(X^n_t)_\#\mu_{\mathrm{in}},
\]
where \(X^n\) is the flow generated by \(V[\mu^n]\).  We claim that, after
choosing a constant \(A_0\) depending only on
\( \M,b,\alpha,\|\mu_{\mathrm{in}}\|_{C^{0,\alpha}}\) and decreasing \(T_1\) if necessary,
\begin{equation}\label{eq:holder-picard-claim}
\|\mu^n\|_{L^\infty([0,T_1];C^{0,\alpha})}
\le A_0
\qquad\forall n\ge0 .
\end{equation}
Assume the claim holds for some \(n\).  By
\Cref{lem:wellposed_kernel_estimates},
\[
\|V[\mu^n_t]\|_{C^{1,\alpha}}
\le
C\bigl(1+\|\mu^n_t\|_{C^{0,\alpha}}\bigr)^2
\le L
\qquad\forall t\in[0,T_1]
\]
for \(L=L( \M,b,\alpha,A_0)\).  Hence
\[
\|X^n_t\|_{C^{1,\alpha}}+\|(X^n_t)^{-1}\|_{C^{1,\alpha}}
\le C\exp(CLT_1)
\qquad\forall t\in[0,T_1].
\]
 Let \(JX_t^n\) denote the Riemannian Jacobian of \(X_t^n\).
The same estimate gives
\[
\left\|\frac1{ JX_t^n}\right\|_{C^{0,\alpha}}
\le C\exp(CLT_1).
\]
In particular, exploiting the representation formula
\[
\mu^{n+1}_t
=
\left(\frac{\mu_{\mathrm{in}}}{ JX_t^n}\right)\circ (X_t^n)^{-1},
\]
we conclude
\[
\begin{aligned}
\|\mu^{n+1}_t\|_{C^{0,\alpha}}
&\le
C\left(1+\|(X_t^n)^{-1}\|_{C^{1,\alpha}}^\alpha\right)
\left\|\frac1{ JX_t^n}\right\|_{C^{0,\alpha}}
\|\mu_{\mathrm{in}}\|_{C^{0,\alpha}}  \\
&\le
C\exp(CLT_1)\|\mu_{\mathrm{in}}\|_{C^{0,\alpha}}
\le A_0,
\end{aligned}
\]
provided that \(A_0\) is fixed large enough and then \(T_1>0\) is chosen
sufficiently small.  This proves
\eqref{eq:holder-picard-claim}.  
For each fixed \(t\), the uniform \(C^{0,\alpha}\) bound makes
\(\{\mu_t^n\}_{n\ge0}\) relatively compact in \(C^0(\M)\) by the
Ascoli-Arzel\'a theorem. Moreover, the convergence established in the
Yudovich construction shows that every uniform subsequential limit of these
Picard iterates coincides with \(\mu_t\).
Hence
\(\mu_t\in C^{0,\alpha}\) and
\(\|\mu_t\|_{C^{0,\alpha}}\le A_0\) for \(t\in[0,T_1]\), by lower
semicontinuity of the Hölder seminorm under uniform convergence.

\smallskip
\noindent\textbf{Step 2.}
Let
\[
T_\ast:=\sup\{s\in[T_1,T]:
\mu\in L^\infty([0,s);C^{0,\alpha})\}.
\]
We prove that \(T_\ast=T\).  All estimates below are first made on
\([0,t]\) with \(t<T_\ast\), and the constants are independent of \(t\).  Let
\(X_t\) be the flow generated by \(V[\mu_t]\), and set
\[
\tilde\mu_t:=\mu_t\circ X_t,\qquad
Z(t):=\int_0^t\|\nabla V[\mu_s]\|_{L^\infty}\,ds .
\]
The flow bounds follow from Gronwall's lemma applied to
\(\frac{d}{dt}\operatorname{dist}_{ \M}(X_t(x),X_t(y))\le
\|\nabla V[\mu_t]\|_{L^\infty}\operatorname{dist}_{ \M}(X_t(x),X_t(y))\), and similarly for the
backward flow:
\[
\operatorname{Lip}(X_t)+\operatorname{Lip}(X_t^{-1})\le C e^{Z(t)},
\]
where the constant only accounts for working in finitely many charts on
 \(\M\). Since \(\mu_t=\tilde\mu_t\circ X_t^{-1}\), this gives
\[
\begin{aligned}
\|\mu_t\|_{C^{0,\alpha}}
&\le \|\tilde\mu_t\|_{L^\infty}
+[\tilde\mu_t]_{C^{0,\alpha}}\operatorname{Lip}(X_t^{-1})^\alpha  \\
&\le C e^{\alpha Z(t)}\|\tilde\mu_t\|_{C^{0,\alpha}} .
\end{aligned}
\]
Along the flow,
\[
\partial_t\tilde\mu_t
=\left(\partial_t\mu_t+V[\mu_t]\cdot\nabla\mu_t\right)\circ X_t
=-\left(\mu_t\divg V[\mu_t]\right)\circ X_t .
\] 
In our
notation \(V[\mu_t]=-\nabla\log u_t\), hence
\(\divg V[\mu_t]=-\Delta\log u_t\), and therefore
\[
\partial_t\tilde\mu_t=\tilde\mu_t\,(\Delta\log u_t)\circ X_t .
\]
Using \eqref{eq:short_chizat_G}, we write this equation as
\begin{equation}
\label{eq:wp_A_B}
\partial_t\tilde\mu_t
=B_t\tilde\mu_t+A_t\tilde\mu_t^2,
\qquad
A_t:=-\frac1{b^2(u_t\circ X_t)},\qquad
B_t:=\frac1{b^2}-|\nabla\log u_t|^2\circ X_t .
\end{equation}
We now only need rough Hölder bounds on the coefficients.  From
\eqref{eq:wellposed_fixed_b_basic}, \(\kappa_b\le u_t\le M_T\) and
\(\|\nabla u_t\|_{L^\infty}\le C\).  Thus \(u_t\) and \(1/u_t\) are bounded in
\(C^{0,\alpha}\).  Moreover, the log-Lipschitz bound for \(\nabla u_t\) in
\Cref{lem:wellposed_kernel_estimates} implies
\([\nabla u_t]_{C^{0,\alpha}}\le C\), since for all $\alpha \in (0,1)$
\(r\log(2+r^{-1})\le Cr^\alpha\) on the compact manifold
 \(\M\).  Hence, we have
\[
\left\|\frac1{u_t}\right\|_{C^{0,\alpha}}
+\left\||\nabla\log u_t|^2\right\|_{C^{0,\alpha}}
\le C .
\]
Composing with \(X_t\) and using \(\operatorname{Lip}(X_t)\le Ce^{Z(t)}\), we get
\[
\|A_t\|_{L^\infty}+\|B_t\|_{L^\infty}\le C,
\qquad
[A_t]_{C^{0,\alpha}}+[B_t]_{C^{0,\alpha}}
\le Ce^{\alpha Z(t)} .
\]
Since \(\|\tilde\mu_t\|_{L^\infty}=\|\mu_t\|_{L^\infty}\le M_T\), the product
estimate in \(C^{0,\alpha}\) gives
\[
\|B_t\tilde\mu_t+A_t\tilde\mu_t^2\|_{C^{0,\alpha}}
\le
C\|\tilde\mu_t\|_{C^{0,\alpha}}+Ce^{\alpha Z(t)} .
\]
Therefore
\[
\frac{d}{dt}\|\tilde\mu_t\|_{C^{0,\alpha}}
\le
C\|\tilde\mu_t\|_{C^{0,\alpha}}+Ce^{\alpha Z(t)} .
\]
Gronwall's lemma gives
\[
\|\tilde\mu_t\|_{C^{0,\alpha}}
\le
C\|\mu_{\mathrm{in}}\|_{C^{0,\alpha}}
+C\int_0^t e^{\alpha Z(r)}\,dr
\le Ce^{\alpha Z(t)},
\]
because \(Z\) is nondecreasing and \(t\le T\). Hence
\[
\|\mu_t\|_{C^{0,\alpha}}
\le Ce^{2\alpha Z(t)} .
\]
On the other hand, by \Cref{lem:wellposed_log_lipschitz},
\[
\|\nabla V[\mu_t]\|_{L^\infty}
\le
C\left(1+\log(2+\|\mu_t\|_{C^{0,\alpha}})\right).
\]
Combining the last two estimates,
\[
Z'(t)=\|\nabla V[\mu_t]\|_{L^\infty}
\le C\bigl(1+Z(t)\bigr).
\]
Another application of Gronwall's lemma gives a finite bound for \(Z\), and
hence for \(\sup_{t\in[0,T_\ast)}\|\mu_t\|_{C^{0,\alpha}}\). If
\(T_\ast<T\), choose \(t_0<T_\ast\) close enough that the short-time
construction of Step 1, restarted from \(\mu_{t_0}\), exists beyond
\(T_\ast\). By uniqueness this construction is \(\mu\), contradicting the
definition of \(T_\ast\), and therefore \(T_\ast=T\).

\smallskip
\noindent\textbf{Step 3.}
We now bootstrap.  Assume, for some \(0\le q<k\), that
\(\mu\in L^\infty([0,T];C^{q,\alpha})\).  By
\Cref{lem:wellposed_kernel_estimates},
\[
u_t=\mathcal K_b\mu_t\in C^{q+2,\alpha},
\qquad
V[\mu_t]=-\nabla\log u_t\in C^{q+1,\alpha},
\]
uniformly in \(t\).  Thus the characteristic flow \(X_t\), and its inverse,
are uniformly bounded in \(C^{q+1,\alpha}\).  For
\(\tilde\mu_t:=\mu_t\circ X_t\), Step 2 gives \eqref{eq:wp_A_B}, i.e.:
\[
\partial_t\tilde\mu_t=B_t\tilde\mu_t+A_t\tilde\mu_t^2,
\qquad
A_t=-\frac1{b^2(u_t\circ X_t)},\qquad
B_t=\frac1{b^2}-|\nabla\log u_t|^2\circ X_t .
\]
Since \(u_t\ge\kappa_b\), the algebra and composition estimates in Hölder
spaces give
\[
\|A_t\|_{C^{q+1,\alpha}}+\|B_t\|_{C^{q+1,\alpha}}\le C,
\qquad
\|\tilde\mu_t\|_{C^{q,\alpha}}\le C .
\]
Consequently a product estimate yields
\[
\|B_t\tilde\mu_t+A_t\tilde\mu_t^2\|_{C^{q+1,\alpha}}
\le C\bigl(1+\|\tilde\mu_t\|_{C^{q+1,\alpha}}\bigr).
\]
The key point is that all terms containing \(q+1\) derivatives of
\(\tilde\mu_t\) are linear in the top norm, while every remaining factor is already
controlled in \(C^{q,\alpha}\).  Hence,
\[
\frac{d}{dt}\|\tilde\mu_t\|_{C^{q+1,\alpha}}
\le C\bigl(1+\|\tilde\mu_t\|_{C^{q+1,\alpha}}\bigr).
\]
Gronwall's lemma gives
\(\tilde\mu\in L^\infty([0,T];C^{q+1,\alpha})\), and composition with
\(X_t^{-1}\) gives
\(\mu\in L^\infty([0,T];C^{q+1,\alpha})\).  Iterating from Step 2 proves the
\(C^{k,\alpha}\) bound.

If \(\mu_{\mathrm{in}}\in C^\infty\), the same argument applies for every \(k\).  The
equation then gives smoothness in time by differentiating
\(\partial_t\mu=-\divg(\mu V[\mu])\) and using the already obtained spatial
regularity, arguing inductively on the number of time derivatives.
\end{proof}

\end{document}
